% Part: normal-modal-logic
% Chapter: completeness
% Section: lindenbaums-lemma

\documentclass[../../../include/open-logic-section]{subfiles}

\begin{document}

\olfileid{nml}{com}{lin}

\olsection{مبرهنة ليندنباوم}

تقرر مبرهنة ليندنباوم أن كل مجموعة $\Sigma$-متسقة من ال!!{formula}s
محتواة في مجموعة $\Sigma$-متسقة \emph{تامة} واحدة على الأقل. وسيبين
بناؤنا للنموذج القانوني أنه تقابل كل مجموعة $\Sigma$-متسقة تامة
$\Delta$ عالمًا في النموذج القانوني تصدق فيه كل ال!!{formula}s في
$\Delta$، ولا تصدق فيه إلا هي. ومن ثم تضمن مبرهنة ليندنباوم صدق كل
مجموعة $\Sigma$-متسقة عند عالم ما في النموذج القانوني.

\begin{thm}[مبرهنة ليندنباوم]\ollabel{thm:lindenbaum}
  إذا كانت~$\Gamma$ $\Sigma$-متسقة، فثمة مجموعة $\Sigma$-متسقة تامة
  $\Delta$ توسّع~$\Gamma$.
\end{thm}

\begin{proof}
لتكن~$!A_0$، $!A_1$، \dots{} قائمة شاملة لجميع صيغ اللغة (ويجوز
التكرار). فمثلًا، نبدأ بإدراج~$\Obj p_0$، وفي كل مرحلة~$n \ge 1$
ندرج الصيغ المنتهية الكثيرة التي لا يتجاوز طولها~$n$ ولا تستعمل إلا
المتغيرات~$\Obj p_0$، \dots،~$\Obj p_n$. ونعرّف مجموعات من
ال!!{formula}s~$\Delta_n$ بالاستقراء على~$n$، ثم نضع
$\Delta = \bigcup_n \Delta_n$. نضع أولًا~$\Delta_0 = \Gamma$.
وبافتراض أن~$\Delta_n$ قد عُرّفت، نعرّف~$\Delta_{n+1}$ هكذا:
\[
\Delta_{n+1} =
\begin{cases}
  \Delta_n \cup \{!A_n\}, & \text{إذا كانت $\Delta_n \cup \{!A_n\}$
    $\Sigma$-متسقة؛} \\
  \Delta_n \cup \{\lnot !A_n\}, & \text{فيما عدا ذلك.}
\end{cases}
\]
والآن لتكن~$\Delta = \bigcup_{n=0}^\infty \Delta_n$.

علينا أن نبين أن هذا التعريف ينتج بالفعل مجموعة~$\Delta$ ذات الخصائص
المطلوبة؛ أي إن~$\Gamma \subseteq \Delta$ وإن~$\Delta$ $\Sigma$-متسقة
تامة.

من الواضح أن~$\Gamma \subseteq \Delta$، لأن~$\Delta_0 \subseteq
\Delta$ بالبناء و$\Delta_0 = \Gamma$. بل إن~$\Delta_n \subseteq
\Delta$ لكل~$n$، لأن~$\Delta$ اتحاد جميع~$\Delta_n$. (ولأننا نضيف في
كل خطوة من البناء !!a{formula} إلى المجموعة المبنية من قبل، يكون
$\Delta_n \subseteq \Delta_{n+1}$؛ ومن ثم، بتعدي~$\subseteq$، يكون
$\Delta_n \subseteq \Delta_m$ متى كان~$n \le m$.) وفي كل مرحلة من
البناء نضيف إما~$!A_n$ أو~$\lnot !A_n$، وتظهر كل !!{formula} مرة واحدة
على الأقل في قائمة جميع~$!A_n$. ومن ثم، لكل~$!A$، إما
$!A \in \Delta$ أو~$\lnot !A \in \Delta$، فتكون~$\Delta$ تامة بحسب
التعريف.

ويبقى أن نبين أن~$\Delta$ $\Sigma$-متسقة. ولهذا نبين أن (أ) عدم اتساق
$\Delta$ مع~$\Sigma$ يقتضي عدم اتساق~$\Delta_n$ ما مع~$\Sigma$، وأن
(ب) كل~$\Delta_n$ $\Sigma$-متسقة.

فافترض أن~$\Delta$ $\Sigma$-غير متسقة. عندئذ
$\Delta \Proves[\Sigma] \lfalse$؛ أي توجد~$!A_1$، \dots،~$!A_k
\in \Delta$ بحيث
$\Sigma \Proves !A_1 \lif (!A_2 \lif \cdots (!A_k \lif
\lfalse)\dots)$. ولأن~$\Delta = \bigcup_{n=0}^\infty \Delta_n$،
يكون كل~$!A_i \in \Delta_{n_i}$ لبعض~$n_i$. ولتكن~$n$ أكبر هذه
الفهارس. لما كان~$n_i \le n$، كان~$\Delta_{n_i} \subseteq \Delta_n$.
وهكذا تنتمي جميع~$!A_i$ إلى~$\Delta_n$ واحدة. وهذا يعني
$\Delta_n \Proves[\Sigma] \lfalse$؛ أي إن~$\Delta_n$
$\Sigma$-غير متسقة.

ولإثبات أن كل~$\Delta_n$ $\Sigma$-متسقة، نستعمل استقراءً بسيطًا
على~$n$. لدينا~$\Delta_0 = \Gamma$، وقد افترضنا أن~$\Gamma$
$\Sigma$-متسقة؛ فتصدق الدعوى عندما~$n = 0$. والآن افترض صدقها عند~$n$،
أي إن~$\Delta_n$ $\Sigma$-متسقة. تكون~$\Delta_{n+1}$ إما
$\Delta_n \cup \{!A_n\}$ إذا كانت هذه $\Sigma$-متسقة، وإلا فهي
$\Delta_n \cup \{\lnot!A_n\}$. وفي الحالة الأولى تكون~$\Delta_{n+1}$
$\Sigma$-متسقة بوضوح. لكن، بحسب
\olref[prf][con]{prop:consistencyfacts}\olref[prf][con]{prop:consistencyfacts-c}،
تكون إما~$\Delta_n \cup \{!A_n\}$ أو~$\Delta_n \cup \{\lnot!A_n\}$
متسقة، ومن ثم تكون~$\Delta_{n+1}$ متسقة في الحالة الأخرى أيضًا.
\end{proof}

\begin{cor}\ollabel{cor:provability-characterization}
  $\Gamma \Proves[\Sigma] !A$ إذا وفقط إذا كان~$!A \in \Delta$ لكل
  مجموعة $\Sigma$-متسقة تامة~$\Delta$ توسع~$\Gamma$ (بما في ذلك حالة
  $\Gamma = \emptyset$، فنحصل عندئذ على توصيف آخر للنسق الجهوي
  $\Sigma$).
\end{cor}

\begin{proof}
  افترض أن~$\Gamma \Proves[\Sigma] !A$، ولتكن~$\Delta$ أي مجموعة
  $\Sigma$-متسقة تامة توسع~$\Gamma$. إذا كان~$!A \notin \Delta$،
  كان~$\lnot!A \in \Delta$ بالقصوية، ومن ثم
  $\Delta \Proves[\Sigma] !A$ (بالرتابة) و
  $\Delta \Proves[\Sigma] \lnot!A$ (بالانعكاسية)، فتكون~$\Delta$ غير
  متسقة. وبالعكس، إذا كان~$\Gamma \Proves/[\Sigma] !A$، كانت
  $\Gamma \cup \{\lnot!A\}$ $\Sigma$-متسقة، وبمبرهنة ليندنباوم توجد
  مجموعة متسقة تامة~$\Delta$ توسع~$\Gamma \cup \{\lnot!A\}$. وبالاتساق،
  $!A \notin \Delta$.
\end{proof}

\end{document}
